\scanpage{143}
\deleted{Ceci \unsure{prouvé|posé}}

D'autre part on a

\[
\unsure{\mathfrak{Z}|\mathfrak{F}|\mathfrak{z}}(G\times G')=\mathfrak{Z}(G)\times\mathfrak{Z}(G')
\]

et \deleted{\unsure{les}} l'ens. des sousgroupes $\in\mathscr C$ d'un $G\in\mathscr C$ satisfait la condition minimale.

Ces propriétés formelles suffisent pour construire un \uline{pro-objet} de $\mathscr C$, \deleted{\unsure{au}} au moyen d'un système projectif \unsure{de épimorphismes|d'épimorphismes}, qu'on \deleted{\unsure{pe}} désignera par $\pi_1^{\mathscr C}(S,\xi)$.

\uline{Exemple I} $\mathscr C={}$ \unsure{cas des groupes|ens. des groupes} définis par \unsure{des|un|les} \unsure{groupes finis ordinaires|groupe fini ordinaire}. On retrouve le groupe fondamental.

\uline{Exple II} $S$ est un préschéma au dessus d'un corps $k$. $\mathscr C$ provient d'une catégorie $\mathscr C_0$ de $k$-groupes finis (par exemple \deleted{\unsure{ét--|tous}} les $k$-groupes finis --- ce qui \deleted{\unsure{est}} \inserted{semble} raisonnable surtout si $k$ sep. clos) stable par produits, noyaux de couples, images.

% archive sheet number: 142

